% secV10_conclusion.tex -- Section 10. CONTAINER (27 Oct 2026): text to be written by the author.
% Word budget: 400 words (JPART count: running text, footnotes and citations as printed).
% Outline, numbers and sources: referee_EE_2026-10/ESQUELETO.md, section 10. The archived text
% (_arquivo_texto_IA_2026-10-27/) is not a source.

\section{Conclusion}\label{sec:conclusion}

S\~ao Paulo's Municipal Education Plan combines a population-coverage requirement with a commitment to serve registered demand, whichever implies the greater obligation. Yet the law leaves the conventions needed to turn those requirements into an observable compliance test unresolved. A subsequent ordinance specified a December 31 monitoring date \citep{sme2016portaria7972} but not a formula. The municipality's 2023 review classified Goal~5 as achieved \citep{saopaulo2023conferencia} and reported that its 75\% criterion had been met since 2015, although that criterion is not equivalent to the statutory requirement.

Across our measurement grid, Goal~5 compliance is specification-dependent in 2019--2022, whereas no specification in that grid indicates compliance in any evaluable year in 2010--2018; before September 2019, only the published waiting-list series is available. Seven of the 12 principal specifications yield compliance in at least one year, and five in none. In the wider grid, none of the 648 evaluable specification-year observations from 2010--2018 qualifies. The review's claim of achievement since 2015 is not supported by any included compliance specification. Its reported figures are largely reconstructible from a December demand-service ratio, but that ratio does not identify or test the binding requirement.

The ambiguity is not solely textual. The first crossing of the review's 75\% threshold moves from 2015 with December counts to 2018 with September counts. Population series determine whether enrolment exceeds the coverage threshold; the law does not select one. The application records, moreover, are incomplete in earlier years and show a shift towards younger ages at first request. Those facts do not alter the compliance exercise based on the published aggregate counts. They do show why a dated waiting-list stock is an incomplete account of access and waiting time.

The new national education plan repeats the demand-based commitment while assigning the development of monitoring indicators to Inep \citep{brasil2026pne}. The case points to a minimum standard for an auditable target: a published formula and population source, explicit definitions of eligible requests and service, consistent enrolment universes, reference and extraction dates, and a reconciliation of entries, exits, and changes in the register. A demand measure should be accompanied by population coverage, because families unable to register are absent from the former. These choices will not remove disagreement about the appropriate standard. They will establish what the reported performance claim means and permit its independent assessment.
